% B8 草稿存档 v1.1（2026-11-02 修订, 依据《报告对比评审与修改意见》D3/D5/M4 + 中文版批注 #10–#13）:
%   [D5] Handling decisions 段压缩: B4 完整数字移出数据章, 指向 Section 6; 落位文本在
%        docs/B8_sec6_footnote_for_D_v1.0.tex（含批注 #12/#13 要求的三方案对照式）。
%   [M4] 类别比记号统一: (≈94:6) → i.e. 6.1% positives (94:6); 图1图注同步。
%   [D3/#10] 出处段轻量压缩: 保留"7,200 条不适用"澄清(批注#10 明确要求保留), 压缩 Quinlan 版说明。
%   [#11] 审计段再收紧: 删两处括号补充(zero anomalies 的展开 / TBG flag 的机制复述), 机制主线不动。
%   未动(护城河): 配对机制、zero anomalies、TBG 100%、B4 结论句、RF 精度保持伏笔。
% 恢复方法:
%   tail -n +14 "docs/B8_sec3b_sick_draft_v1.0.tex" > overleaf/sections/sec3b_sick.tex
% 数字来源: results/stats/sick_missing_audit.csv(B3) / results/b4_missing_comparison.csv(B4) /
%          results/sick_*_summary.csv(主矩阵, D4 已核验) / c_config.py(冻结口径)。
% 交叉引用依赖: \ref{sec:results} 依赖 D 的 sec6 label(见 docs/D7_sec6_results_draft_v1.0.tex)。

\subsection{Sick Dataset}
\label{sec:sick-data}
% B owns this subsection: UCI provenance, label meaning, missingness, and preprocessing.

The \textit{sick} dataset belongs to the UCI Thyroid Disease repository
(ID 102, DOI 10.24432/C5D010), with original data donated by the Garavan
Institute (1986) and redistributed under CC BY 4.0; we use the Weka-packed
variant (accepted by the course brief). The ``7,200 instances / 5
features'' on the UCI landing page describes the whole thyroid
repository, not this subset. The actual dataset contains
\textbf{3,772 instances $\times$ 30 attributes} (23 nominal, 7 numeric,
including the class). The class is severely imbalanced: 3,541
\textit{negative} versus 231 \textit{sick-euthyroid}, i.e.\ 6.1\%
positives (94:6); we treat \textit{sick} as the positive class
throughout. Two semantic notes we keep strict in wording:
\textit{sick-euthyroid} denotes euthyroid illness --- it is
\emph{not} hypothyroidism --- and \textit{negative} means no thyroid
abnormality detected in this screening, not ``perfectly healthy''. 61
duplicate rows exist in the official distribution; we keep them to remain
consistent with the official instance count and disclose them as a
data-quality issue.

\subsubsection{Missing-value audit}
In total 6,064 cells (5.4\%) are missing, coded as ``?''. The audit
revealed a clean \emph{measurement mechanism}: every laboratory test
appears as a pair --- a \texttt{*\_measured} flag and the value itself ---
and whenever the flag reads \texttt{f} (test not performed), the paired
value is missing \emph{by design}. Across all six lab pairs the
consistency check found \textbf{zero anomalies}: the data is incomplete
but internally consistent. The extreme case is \textbf{TBG, 100\% missing
in all 3,772 rows} (its flag reads \texttt{f} in every row). After
removing TBG and its constant flag, 1,129 rows
(29.9\%) still carry at least one missing value
(Figs.~\ref{fig:sick-dist}--\ref{fig:sick-missing}).

\subsubsection{Handling decisions}
The handling follows the mechanism. (i)~TBG and its constant flag are
\emph{removed}, not imputed --- a column with zero observations cannot be
imputed without fabricating data; column removal involves no statistics
and is leakage-free. (ii)~All remaining columns are imputed
\emph{inside} the imblearn pipeline (numeric by mean, nominal by mode),
so imputation statistics are fitted on training folds only, alongside
encoding, scaling and SMOTE; the whole project evaluates on group-locked
splits. We also compared three ways of handling the remaining gaps,
using the decision tree as the test bed: imputing inside the folds,
dropping incomplete rows, and treating ``missing'' as its own category
(scikit-learn trees cannot accept \texttt{NaN} values, unlike Weka's
J48). Dropping rows turned out to discard 29.9\% of the patients for no
measurable benefit, so we impute. The full comparison, with per-fold
numbers and the AUC contrast, is reported together with the results in
Section~\ref{sec:results}.

Finally, among the four base classifiers the random forest is the only
one whose precision is preserved under SMOTE (0.94$\to$0.91, vs.\
0.95$\to$0.60 for a single tree); Section~\ref{sec:results} sharpens this
account with the ensemble study, where AdaBoost shares the property and
plain bagging does not.

\begin{figure}[htbp]
\centering
\begin{minipage}{0.48\linewidth}
  \centering
  \includegraphics[width=\linewidth]{figures/fig_b_class.pdf}
  \caption{Class distribution of sick (3,541:231, 6.1\% positives; cv10,
  positive = sick-euthyroid).}\label{fig:sick-dist}
\end{minipage}\hfill
\begin{minipage}{0.48\linewidth}
  \centering
  \includegraphics[width=\linewidth]{figures/fig_b_missing.pdf}
  \caption{Per-column missing rate of sick; TBG (red) is 100\% missing
  and removed rather than imputed.}\label{fig:sick-missing}
\end{minipage}
\end{figure}

\begin{figure}[htbp]
\centering
\includegraphics[width=0.62\linewidth]{figures/fig_b_box.pdf}
\caption{TSH (log scale) and FTI by class (raw data): heavy class overlap
(median TSH 1.4 vs.\ 1.5; FTI 107 vs.\ 108) shows no single lab value
separates the classes, motivating multivariate models.}\label{fig:sick-box}
\end{figure}
